\section{模板使用说明}

本模板依据广西大学学位论文格式规范，基于 \TeX{} Live 2025 制作（最低测试至 \TeX{} Live 2023），默认行距对应 Word 中的 1.25 倍行距效果\footnote{"文档网格——只指定行网格——每页 44 行"下的表现。}。如需切换至 25 磅行距，可在正文开始前调用 \verb|\gxustretchwide| 命令；如需切换回默认行距，调用 \verb|\gxustretchnormal|。需要注意的是，25 磅行距仅调整正文文字的行距，其余格式参数未做相应匹配，实际效果可能欠佳，不建议使用。

\subsection{声明}

本模板作者并非广西大学在读或毕业学生，与广西大学不存在任何隶属关系。本模板参考了学校内流传的不知名 \LaTeX{} 模板，并在此基础上进行了大量修改与重构；若原作者认为本模板侵犯了其权益，请联系作者，将立即删除相关内容。

模板中使用的广西大学题字图片提取自广西大学官方发布的 Word 学位论文模板，仅供本模板排版使用，\textbf{严禁用于其他任何用途}。作者无意侵犯广西大学的任何权益，若校方认为此使用方式不当，请联系作者处理。

由于作者与广西大学并无关联，无法持续获取最新格式规范，\textbf{本模板在发布后将不再主动维护}。若广西大学日后更新学位论文格式要求，本模板可能无法及时跟进。欢迎有意愿的同学或开发者在本模板基础上继续维护与更新，\textbf{更新时请在显著位置保留对原作者的署名：薛定谔之花/Schrodinger Blume}。

\subsection{字体}

本模板使用中易系列字体排版，包括中易宋体、中易黑体与中易隶书。请注意，上述字体均为微软公司的商业字体，受版权保护，\textbf{请勿将字体文件用于商业用途，亦请勿将其上传至公开网络}。

若您的操作系统为 Windows，上述字体通常已随系统预装，模板将自动识别并加载，无需额外操作。

若您的操作系统为 macOS 或 Linux，则需自行获取字体并安装。若您使用在线编译平台（如 Overleaf、TeXPage\footnote{TeXPage 平台支持中易宋体及中易黑体，仅需上传中易隶书字体文件。}等），则应将字体文件放置于项目根目录下的 \texttt{字体/} 文件夹中。文件命名规则如下：
\begin{itemize}
    \item 中易宋体：\texttt{simsun.ttc}（优先）或 \texttt{simsun.ttf}
    \item 中易黑体：\texttt{simhei.ttf}
    \item 中易隶书：\texttt{simli.ttf}
\end{itemize}

注意文件名及扩展名须\textbf{全部小写}，否则将无法被正确识别。若上述字体均无法获取，模板将自动回退至系统默认字体，仍可正常编译，但排版效果或与预期有所差异\footnote{宋体和黑体理论上将退回至 Fandol 系列字体，隶书将用加粗黑体代替。}。

\subsection{标题的使用}

本模板支持四级正式编号标题（\verb|\section| 至 \verb|\paragraph|），标题文字与编号之间的间距默认为一个汉字宽度（\verb|1em|）。如需调整为半个汉字宽度，可在正文开始前调用 \verb|\gxuafternamenarrrow| 命令。

\subsubsection{五级及以上标题的处理}

根据广西大学模板要求，四级标题（\verb|\paragraph|）之后，应使用数字加圆括号或字母标注五级及以上层次。本模板为此预定义了以下有序列表样式，可通过 \verb|\begin{enumerate}[样式名]| 调用：

\begin{itemize}
    \item \verb|[kuohao]|（括号）：用于第五级，编号格式为（1）（2）（3）……
    \item \verb|[daxie]|（大写字母）：用于第六级，编号格式为 A、B、C……
    \item \verb|[xiaoxie]|（小写字母）：用于第七级，编号格式为 a、b、c……
    \item \verb|[shuzi]|（数字带句点）：备用数字样式，与默认样式的区别在于换行后不缩进
\end{itemize}

以下展示各样式的实际效果：

\paragraph{四级标题}

五级标题使用 \verb|[kuohao]| 样式：
\begin{enumerate}[kuohao]
    \item 五级标题示例
    \item 五级标题示例
\end{enumerate}

六级标题使用 \verb|[daxie]| 样式：
\begin{enumerate}[daxie]
    \item 六级标题示例
    \item 六级标题示例
\end{enumerate}

七级标题使用 \verb|[xiaoxie]| 样式：
\begin{enumerate}[xiaoxie]
    \item 七级标题示例
    \item 七级标题示例
\end{enumerate}

\verb|[shuzi]| 样式与默认数字样式的换行缩进对比：
\begin{enumerate}[shuzi]
    \item 数字样式（无缩进）：数字数字数字数字数字数字数字数字数字数字数字数字数字数字数字数字数字数字数字数字数字数字数字数字数字
    \item 数字样式（无缩进）：数字数字数字数字数字数字数字数字数字数字数字数字数字数字数字数字数字数字数字数字数字数字数字数字数字
\end{enumerate}
\begin{enumerate}
    \item 默认样式（有缩进）：数字数字数字数字数字数字数字数字数字数字数字数字数字数字数字数字数字数字数字数字数字数字数字数字数字
    \item 默认样式（有缩进）：数字数字数字数字数字数字数字数字数字数字数字数字数字数字数字数字数字数字数字数字数字数字数字数字数字
\end{enumerate}